% Page budget: 1.55 pages | Owns: negation, OBC derivation, additional-state scope, implementation, and true-parameter condition.
% WRITING GUIDE
% Purpose: Explain negation, derive the final state-level orthogonal-by-construction map, and delimit its interpretability guarantee.
% Start here: Decisions D-185, D-190, D-192, D-194, D-198, D-202, D-203, D-222, D-229, and D-232 in docs/decisions.md.
% Supporting material: Current state-level OBC derivations under scripts/gantry/orthogonal-by-construction/documentation and their thesis-readiness guidance.
% Mandatory papers: Read Györök's projection-regularization paper and orthogonal-by-construction paper in full before writing; keep their guarantees and assumptions separate.
% Research-plan anchor: Preserve Aspect 3's goal of protecting physical interpretability; state clearly how the final construction differs from the proposed regularisation.
% Current decisions: Reduced identifiable coordinates, data-derived reference tuples, a tangent-only one-step physical-state projection, per-epoch refresh, and no affine comparison arm.
% Choices to justify: Protected parameter coordinates, reference-set construction, projected rows, rank tolerance, exclusion of the affine offset, refresh frequency, and treatment of additional states.
% Implementation audit: Read model_augmentation/fit_systems/obc.py, scripts/gantry/gantry_dynamic/obc_gantry.py, obc_diagnostics.py, and the OBC attachment in model.py before fixing the derivation.
% Reference check: Verify the original projection method and separate its claims from the state-level, nonlinear-parameter, LPV, and additional-state extensions developed here.
% Cautions: docs/orthogonal-projection-plan.md describes an older soft-penalty route; one-step state orthogonality is not trajectory-output orthogonality or guaranteed true recovery.
% Planned figures: New geometric projection figure with a physical and additional-state partition inset; show the remaining uniqueness limitation as clearly as the protected direction.
% Section is finished when: The protected space, coordinate frame, gradient treatment, added-state scope, and true-recovery condition are explicit.
\section{Interpretability by Orthogonal Construction}\label{sec:obc}

The split between $f_{\mathrm{base}}$ and $f_{\mathrm{aug}}$ in
\eqref{eq:aug_transition} is not unique (Section~\ref{sec:aug_structure}), so
the estimate of $\theta_{\mathrm{base}}$ need not keep its physical meaning.
We constrain the split by adapting the orthogonal-by-construction (OBC)
parametrisation of Gy\"or\"ok et al.~\cite{gyorok2026obc} from input-output
models linear in their parameters to the self-scheduled state equation.
Section~\ref{sec:obc_negation} identifies the directions the network can
negate, Section~\ref{sec:obc_map} constructs the correction, and
Section~\ref{sec:obc_scope} states what the construction guarantees and when
it recovers the true parameters.

\subsection{Negation in the Augmented State Equation}\label{sec:obc_negation}
To first order, a change of $\theta_{\mathrm{base}}$ moves the one-step
transition along the columns of the transition sensitivity
\begin{equation}
 \Phi_{\bar\theta}(\tilde x,u)
 =\left.\frac{\partial f_{\mathrm{base}}(\theta_{\mathrm{base}},\tilde x,u)}
   {\partial\theta_{\mathrm{base}}}\right|_{\theta_{\mathrm{base}}=\bar\theta_{\mathrm{base}}}
   \in\R^{6\times10},
\label{eq:obc_sens}
\end{equation}
where $\bar\theta_{\mathrm{base}}$ is the expansion point, as
in~\cite[Eq.~(15)]{gyorok2025l4dc}.
A learned write can cancel a parameter displacement $\delta\theta$ along these
columns:
\begin{equation}
\begin{aligned}
 &f_{\mathrm{base}}(\bar\theta_{\mathrm{base}}+\delta\theta,\tilde x,u)
 +\bigl[f_{\mathrm{aug}}(\theta_{\mathrm{aug}},\hat x,u)
 -\Phi_{\bar\theta}(\tilde x,u)\,\delta\theta\bigr]\\
 &\quad\approx f_{\mathrm{base}}(\bar\theta_{\mathrm{base}},\tilde x,u)
 +f_{\mathrm{aug}}(\theta_{\mathrm{aug}},\hat x,u).
\end{aligned}
\label{eq:obc_negation}
\end{equation}
The data cannot attribute such a direction to either component, so the
physical estimate depends on the initial parameter
pair~\cite[Ex.~2]{gyorok2026obc}, \cite[Ex.~1]{gyorok2025l4dc}.
We call this negation.
We seek a split whose physical estimate does not depend on that pair.

For the gantry, differentiating \eqref{eq:eom} with respect to the $j$-th
combination gives the acceleration sensitivity
\begin{equation}
 \frac{\partial\ddot q}{\partial\theta_{\mathrm{base},j}}
 =-M(Y)^{-1}\Bigl[\frac{\partial M(Y)}{\partial\theta_{\mathrm{base},j}}\,\ddot q
 +\frac{\partial C}{\partial\theta_{\mathrm{base},j}}\,\dot q
 +\frac{\partial K}{\partial\theta_{\mathrm{base},j}}\,q\Bigr],
\label{eq:obc_gantry_sens}
\end{equation}
which the RK4 step propagates into $\Phi_{\bar\theta}$.
Combinations in $M(Y)$ act through $\ddot q$, those in $C$ through $\dot q$,
and $k_{b,\Sigma}$ through $q$.
A negating write therefore acts as an added inertia, damping or stiffness
force.
The factor $M(Y)^{-1}$ makes these directions depend on the payload position.

The protected directions are those of the ten combinations
$\theta_{\mathrm{base}}$ of Section~\ref{sec:params}, not of the fourteen raw
parameters.
The raw parameters have four structural null directions, so their stacked
sensitivity is rank deficient for any data.
That violates the rank condition~\cite[Assumption~1]{gyorok2026obc} on which
the uniqueness of the physical estimate rests.
Training uses the logarithmic and bounded coordinates of
Section~\ref{sec:params}.
Their map to $\theta_{\mathrm{base}}$ has an invertible Jacobian, so the
sensitivity with respect to them has the same range.
The correction below depends only on that range and is therefore unchanged.

\subsection{Orthogonal-by-Construction Correction}\label{sec:obc_map}
For a baseline linear in its parameters, Gy\"or\"ok et al.\ subtract from the
network output its least-squares fit on the baseline
regressor~\cite[Eqs.~(8), (9), Lemma~3]{gyorok2026obc}:
\begin{equation}
\begin{aligned}
 \hat Y&=\Phi\theta_{\mathrm{base}}+F-\Phi\theta_{\mathrm{aux}},\qquad
 \theta_{\mathrm{aux}}=(\Phi^\top\Phi)^{-1}\Phi^\top F,\\
 0&=\Phi^\top\bigl(F-\Phi\theta_{\mathrm{aux}}\bigr),
\end{aligned}
\label{eq:obc_gyorok}
\end{equation}
where $\Phi$ stacks the regressor and $F$ the network output over the
estimation data.
The corrected network output has no component that a change of
$\theta_{\mathrm{base}}$ could produce.
The orthogonality penalty of~\cite{gyorok2025l4dc} pursues the same aim
through a weight that trades orthogonality against fit.
The construction needs no such weight.

We apply \eqref{eq:obc_gyorok} to the six physical rows of the one-step
transition, which $f_{\mathrm{aug}}$ writes (Section~\ref{sec:aug_structure}).
The baseline has no linear-in-parameters form, because $f_{\mathrm{base}}$ is
rational in $\theta_{\mathrm{base}}$ through $M(Y)^{-1}$.
The regressor is therefore the sensitivity \eqref{eq:obc_sens}, the linear
term of the first-order expansion of~\cite[Eq.~(15)]{gyorok2025l4dc}.
That expansion also has an affine offset, which~\cite[Eq.~(18)]{gyorok2025l4dc}
adds to the protected set.
We do not protect it, because it multiplies no estimated parameter and does
not enter the error of the physical estimate (Section~\ref{sec:obc_scope}).

The regressor is evaluated at states, but only the positions are measured.
Gy\"or\"ok et al.\ evaluate it on measured data~\cite{gyorok2026obc}, and
\cite[Sec.~3]{gyorok2025l4dc} replaces unmeasured states by a simulation of
the baseline.
Such a simulation drifts on the free axes $X$ and $Y$
(Section~\ref{sec:aug_closed_loop}).
A coefficient fitted on the training rollout instead would make the
correction at step $k$ depend on later states of the same window.
We therefore reconstruct a fixed reference tuple
$(\tilde x_i,\bar x_i,u_i)$ from every training sample,
\begin{equation}
 \tilde x_i=\col(q_i,\dot q_i),\qquad q_i=P^{-\top}y_i,\qquad
 \bar x_i=0,\qquad i=1,\dots,N,
\label{eq:obc_tuples}
\end{equation}
where $\dot q_i$ is the fourth-order central difference of $q$ within one
record and $u_i$ is the recorded actuator force.
The tuples are formed in physical units and normalised as the signals.
Each tuple carries its measured $Y_i$, so the $Y$-dependent directions are
sampled over the payload range of the training records.
The recorded input already contains the feedback action, so the controller
does not enter.
Gy\"or\"ok et al.\ permit such an auxiliary evaluation
set~\cite[Sec.~3.2]{gyorok2026obc}.

On the tuples, the construction becomes
\begin{equation}
\begin{aligned}
 \Phi_{\bar\theta}(\tilde X,U)
 &=\col\bigl(\Phi_{\bar\theta}(\tilde x_1,u_1),\dots,
   \Phi_{\bar\theta}(\tilde x_N,u_N)\bigr)\in\R^{6N\times10},\\
 F_{\mathrm{aug}}
 &=\col\bigl(f_{\mathrm{aug}}(\theta_{\mathrm{aug}},\hat x_1,u_1),\dots,
   f_{\mathrm{aug}}(\theta_{\mathrm{aug}},\hat x_N,u_N)\bigr),\\
 \theta_{\mathrm{aux}}&=\Phi_{\bar\theta}(\tilde X,U)^{+}F_{\mathrm{aug}},\\
 0&=\Phi_{\bar\theta}(\tilde X,U)^\top
   \bigl(F_{\mathrm{aug}}-\Phi_{\bar\theta}(\tilde X,U)\,\theta_{\mathrm{aux}}\bigr),\\
 0&=\Phi_{\bar\theta}(\tilde X,U)^\top\,
   \frac{\partial\bigl(F_{\mathrm{aug}}-\Phi_{\bar\theta}(\tilde X,U)\,
   \theta_{\mathrm{aux}}\bigr)}{\partial\theta_{\mathrm{aug}}},
\end{aligned}
\label{eq:obc_ours}
\end{equation}
where $\hat x_i=\col(\tilde x_i,\bar x_i)$ and ${}^{+}$ denotes the
pseudoinverse.
The last line keeps every gradient step of $\theta_{\mathrm{aug}}$ out of the
protected range, which is the property behind~\cite[Eq.~(39)]{gyorok2026obc}.
The fourth line is~\cite[Lemma~3]{gyorok2026obc} on the reference stack and
holds at any rank.
A unique physical estimate needs full column
rank~\cite[Assumption~1]{gyorok2026obc}, which is checked at every rebuild.
Section~\ref{sec:results} reports the rank and condition number on the
training data.

The coefficient follows the network, while the stack is held fixed within an
epoch.
$\theta_{\mathrm{aux}}$ is recomputed at every objective evaluation and
differentiated through.
The stack does not depend on $\theta_{\mathrm{aug}}$, which gives the last
line of \eqref{eq:obc_ours}.
A detached coefficient would lose that line.
Gy\"or\"ok et al.\ update the expansion point at every objective evaluation or
fix it at the nominal parameters~\cite[Sec.~4]{gyorok2025l4dc}.
Updating rebuilds the full stack at every evaluation.
A fixed point does not apply here, because training starts from detuned
parameters (Section~\ref{sec:setup}).
We therefore rebuild $\Phi_{\bar\theta}(\tilde X,U)$ at the current
$\theta_{\mathrm{base}}$ between epochs.
The stack and $\bar\theta_{\mathrm{base}}$ carry no gradient, so
$\theta_{\mathrm{base}}$ is updated only through $f_{\mathrm{base}}$.

The corrected augmented model, with the output map of
\eqref{eq:aug_transition}, is
\begin{equation}
\begin{aligned}
 \begin{bmatrix}\tilde x_{k+1}\\ \bar x_{k+1}\end{bmatrix}
 &=\begin{bmatrix}f_{\mathrm{base}}(\theta_{\mathrm{base}},\tilde x_k,\hat u_k)\\ 0\end{bmatrix}\\
 &\quad+\begin{bmatrix}f_{\mathrm{aug}}(\theta_{\mathrm{aug}},\hat x_k,\hat u_k)
   -\Phi_{\bar\theta}(\tilde x_k,\hat u_k)\,\theta_{\mathrm{aux}}\\
   g_{\mathrm{aug}}(\theta_{\mathrm{aug}},\hat x_k,\hat u_k)\end{bmatrix},
\end{aligned}
\label{eq:obc_model}
\end{equation}
where $\theta_{\mathrm{aux}}$ is the coefficient of the current objective
evaluation.
The correction is evaluated at the rollout point with the frozen expansion
point, as in~\cite[Eq.~(13)]{gyorok2026obc}.
The baseline sensitivity of the additional states is zero, so
$g_{\mathrm{aug}}$ is not corrected.
For prediction, the stack is rebuilt at the trained $\theta_{\mathrm{base}}$.
$\theta_{\mathrm{aux}}$ is then computed once from the trained network and
held fixed.

\subsection{Scope and Condition for True Parameter Recovery}\label{sec:obc_scope}
The guarantee is the orthogonality in \eqref{eq:obc_ours}.
It holds for the one-step stack over the tuples, not at each tuple.
It is local in $\theta_{\mathrm{base}}$, because $\Phi_{\bar\theta}$ is the
tangent at $\bar\theta_{\mathrm{base}}$.
Its inner product is Euclidean in normalised state coordinates, so the state
normalisation sets its metric.

The additional states and the output trajectory lie outside the guarantee.
The tuples lie on $\bar x=0$, so $\theta_{\mathrm{aux}}$ does not see the part
of $f_{\mathrm{aug}}$ that acts through nonzero additional states.
The rows $g_{\mathrm{aug}}$ are never corrected.
They reach the physical rows one step later through $f_{\mathrm{aug}}$.
The output over a window accumulates these contributions through the
recursion and the controller.
One-step state orthogonality therefore does not imply orthogonality of the
output trajectory.
The consistency and zero-covariance results
of~\cite[Thms.~16, 18]{gyorok2026obc} do not transfer to the trained model.

A unique split does not imply recovery of the true parameters.
Let $\theta^*_{\mathrm{base}}$ be the true parameters and $\Delta^*$ the true
one-step discrepancy on the tuples,
\begin{equation}
 \Delta^*=\col\bigl(\tilde x_{i+1}-f_{\mathrm{base}}(\theta^*_{\mathrm{base}},
 \tilde x_i,u_i)\bigr)_{i=1}^{N},
\label{eq:obc_delta}
\end{equation}
where the recorded successor $\tilde x_{i+1}$ is reconstructed with the same
stencil and normalised as the tuples.
Assume noise-free data and $f_{\mathrm{base}}$ affine in
$\theta_{\mathrm{base}}$ near $\theta^*_{\mathrm{base}}$.
Let $\Phi$ be the stack at
$\bar\theta_{\mathrm{base}}=\theta^*_{\mathrm{base}}$, with full column rank.
Assume the corrected one-step model reproduces every recorded successor.
The exact fit and its product with $\Phi^\top$, using \eqref{eq:obc_ours},
give
\begin{equation}
\begin{aligned}
 \Phi\bigl(\hat\theta_{\mathrm{base}}-\theta^*_{\mathrm{base}}\bigr)
 +\widetilde F_{\mathrm{aug}}&=\Delta^*,\\
 \hat\theta_{\mathrm{base}}-\theta^*_{\mathrm{base}}=\Phi^{+}\Delta^*,
 \qquad
 \hat\theta_{\mathrm{base}}=\theta^*_{\mathrm{base}}
 \;&\Longleftrightarrow\;
 \Phi^\top\Delta^*=0,
\end{aligned}
\label{eq:obc_bias}
\end{equation}
where $\widetilde F_{\mathrm{aug}}=F_{\mathrm{aug}}-\Phi\theta_{\mathrm{aux}}$
is the corrected write.
This is the state-level counterpart
of~\cite[Thm.~7, Eq.~(21)]{gyorok2026obc}.
The objective \eqref{eq:aug_loss} is a closed-loop multistep output error, so
\eqref{eq:obc_bias} is a one-step surrogate, not the bias of the trained model.
The condition $\Phi^\top\Delta^*=0$ is a sum over the tuples.
Contributions that are not orthogonal at a single tuple must therefore cancel
over the records.
When the structure of the discrepancy is known, the excitation can be
designed to achieve this~\cite[Remark~5]{gyorok2026obc}.

One overlap measure tests both the recovery condition and the construction.
For a stacked vector $v$ over the tuples and a stack $B$,
\begin{equation}
 \rho(v;B)=\frac{\norm{BB^{+}v}}{\norm{v}},\qquad
 B\in\bigl\{\Phi,\;\begin{bmatrix}\Phi & F_{\mathrm{base}}\end{bmatrix}\bigr\},
\label{eq:obc_overlap}
\end{equation}
where $F_{\mathrm{base}}=\col\bigl(f_{\mathrm{base}}(\bar\theta_{\mathrm{base}},
\tilde x_i,u_i)\bigr)_{i=1}^{N}$ and $\Phi$ is the stack at
$\bar\theta_{\mathrm{base}}$.
The second stack spans the affine set of~\cite[Eq.~(18)]{gyorok2025l4dc}.
For $v=\Delta^*$ and $B=\Phi$, $\rho$ is zero exactly when the recovery
condition holds.
It needs $\theta^*_{\mathrm{base}}$ and is therefore available in simulation
only.
For $v=F_{\mathrm{aug}}$ and $B=\Phi$, $\rho$ is the fraction of the learned
write that the correction removes.
The increase of $\rho$ from $B=\Phi$ to the extended stack measures the
overlap with the nominal baseline response, which the construction leaves
unprotected.
Section~\ref{sec:results} tests whether the split is independent of the
start and, where the condition holds, whether the true parameters are
recovered.
